\begin{table}[htbp]
\centering
\small
\caption{Viés de nomes de atributo (openai/gpt-4o-mini): fidelidade, direção de $\hat{W}$ e consistência.}
\label{tab:bloco1_vies_nomes_atributo_gpt4mini}
\begin{tabular}{llrrrrrrr}
\toprule
Variante & Regime & Fid.\ A & Acur.\ LLM vs real A & $w_1/w_2$ & Consist.\ B & $\kappa$ B & Consist.\ C & $\kappa$ C \\
\midrule
altura / peso & zero-shot & 79,1 (70,7–84,7) & 20,9 & 0,11 & 57,7 & 0,139 & 91,7 & 0,833 \\
feature\_1 / feature\_2 & zero-shot & 78,2 (68,7–84,0) & 31,8 & 0,16 & 81,3 & 0,581 & 83,3 & 0,653 \\
\textbf{x1 / x2} & \textbf{zero-shot} & \textbf{81,3 (78,0–85,3)} & \textbf{32,4} & \textbf{0,09} & \textbf{70,7} & \textbf{0,348} & \textbf{88,0} & \textbf{0,745} \\
\bottomrule
\multicolumn{9}{p{0.95\linewidth}}{\footnotesize Problema A e transferência para B e C por nomes de atributo; os demais fatores ficam na configuração base. Médias das sementes (fidelidade com mín–máx). Fid.\ A = fidelidade (vs LLM) no A; Acur.\ LLM vs real A = acurácia (vs rótulo real) do LLM no A, contexto; $w_1/w_2$ = razão da métrica estimada (Perceptron). Fid.\ A, Acur.\ LLM vs real A e $w_1/w_2$ vêm da Fase A, sempre zero-shot. $\kappa$ = Cohen's $\kappa$ (faixas de Landis e Koch, 1977). Em negrito: a configuração de referência. Em zero-shot (e com nomes de classe neutros) a correspondência entre nome e classe é arbitrária: a acurácia (vs rótulo real) do LLM pode sair abaixo de 50\%; o valor invariante à permutação é max($a$, 1 $-$ $a$).} \\
\end{tabular}
\end{table}
